\section{Conclusion}
\label{sec:conclusion}

We presented a systematic multi-level evaluation of MDL-to-UsdPreviewSurface
material simplification in NVIDIA Isaac Sim, quantifying trade-offs across
four complementary dimensions: pixel-level image quality, rendering performance,
feature-level semantic preservation, and downstream AI task performance.

Our pixel-level analysis demonstrates that UsdPreviewSurface conversions
achieve high fidelity relative to the original MDL renderings, with a mean
PSNR of 35.52\,dB, SSIM of 0.990, and LPIPS of 0.020 across 24 matched image
pairs. These results indicate that the four primary material channels ---
diffuse albedo, roughness, metallic, and normal --- capture the dominant visual
appearance of the tested assets, with degradation concentrated in scenes
featuring complex multi-layer MDL effects such as procedural noise textures and
specular coatings.

Rendering performance measurements reveal that material simplification does not
yield significant throughput or memory benefits for single-object scenes under
the RTX path tracer, with both material types achieving approximately 119\,FPS
and comparable GPU memory consumption. This finding suggests that the primary
motivation for conversion lies in interoperability and portability rather than
rendering efficiency, at least for scenes of the scale tested.

Feature-level evaluation confirms that semantic content is well preserved:
CLIP embeddings achieve a mean cosine similarity of 0.925 across matched pairs,
while DINOv2 --- which is more sensitive to fine-grained texture differences ---
yields 0.872. Zero-shot YOLOv8 detection experiments show that both material
versions trigger comparable detection patterns (MDL: 4/24 images detected, mean
confidence 0.560; UsdPreviewSurface: 7/24, mean confidence 0.513), confirming
that the conversion does not introduce artifacts that degrade object
detectability. These higher-level metrics are particularly important for
practitioners generating synthetic training data, as they directly probe
whether vision models perceive the simplified renderings as equivalent to the
originals.

Based on our findings, we offer the following practical guidelines for
synthetic data practitioners:
\begin{itemize}
    \item For assets with standard PBR materials (diffuse, roughness, metallic,
    normal maps), MDL-to-UsdPreviewSurface conversion is safe and introduces
    negligible visual degradation (PSNR $>$ 35\,dB, SSIM $>$ 0.99).
    \item Assets with complex multi-layer MDL effects should be validated
    individually, as conversion quality varies with material complexity (PSNR
    range: 23--44\,dB across viewpoints).
    \item The conversion enables broader toolchain compatibility --- including
    standard USD viewers, web-based visualization, and GLB export --- without
    sacrificing visual quality for most use cases.
\end{itemize}

\paragraph{Limitations.}
Our evaluation is limited to four chest-of-drawers assets from the Isaac Sim
library, all belonging to the same furniture category. While these assets span
a range of material complexity, further validation on diverse object categories
(e.g., transparent or translucent materials, highly specular metals, organic
surfaces) is needed to establish broader generalizability. Additionally, our
performance benchmarks use single-object scenes; the efficiency impact of
material simplification in large-scale multi-object environments remains an
open question.

\paragraph{Future Work.}
We plan to extend this evaluation along several axes: (1) scaling to larger
asset collections spanning diverse object categories and material types;
(2) measuring the impact on downstream task performance with fine-tuned
detectors and segmentation models trained on simplified vs.\ original
renderings; (3) investigating the cumulative effect of material simplification
combined with mesh simplification and texture compression in the full
USD-to-GLB export pipeline; and (4) benchmarking rendering performance in
complex multi-object scenes where material count may become a bottleneck.
